\section*{Preamble}

In the previous chapter, we have seen that \acsp{TEE} present some operational
challenges that need to be properly addressed when deploying workloads on
untrusted third-party infrastructures such as public clouds. As a follow-up,
this work highlights a specific and non-obvious limitation of \acs{TEE}
architectures that, if overlooked, might subvert the security guarantees
provided by enclaves.

This research is the result of several discussions with Fritz Alder during the
design of the self-revocation scheme that will be presented in \cref{ch:v2x}. As
will be more clear later, our design required computing timeouts and other
operations based on the current time, which initially did not seem to be a
problem but then, upon closer investigation, we stumbled upon some issues: How
can we provide a reliable representation of elapsed time to an enclave, and how
can we synchronize time among a network of enclaves? Indeed, if we look at the
\acs{TEE} threat model, the infrastructure where enclaves run is untrusted,
which means that privileged software layers such as the host operating system
and the hypervisor are assumed to be either malicious or compromised. Hence,
fetching the time from the untrusted environment is dangerous and should be
avoided in favor of more reliable sources. However, some \acsp{TEE} do not
natively provide access to a trusted time source and, even when they do, it is
difficult to safely provide time values to enclaves without being altered by an
attacker, thus making such values partially or completely lose their meaning.

Therefore, measuring time inside an enclave is a delicate task that needs to
take into account factors such as the \acs{TEE} architecture being used, the
time source from which time is read, and the specific use case considered. By
contrast, however, the current trend in confidential computing is to move
existing workloads to \acsp{TEE} with a ``lift-and-shift'' approach, i.e.,
without making any modifications to application code. This trend has been more
pronounced in recent years, with the advent of \acs{VM}-based \acsp{TEE} for
cloud workloads (cf.~\cref{ch:cloud}). Unfortunately, this approach is not
always a good idea as it might provide a false sense of security to application
owners: Considering the usage of time, if an application uses timestamps for
some security-critical operations (e.g., as an entropy source for creating a
cryptographic key), then an attacker that has control over the time source could
potentially exploit the application by providing it with carefully-chosen
values. For instance, most libraries and frameworks for Intel \acs{SGX} rely on
the untrusted time provided by the operating system, thus applications running
in \acs{SGX} enclaves are susceptible to this class of attacks by default and
might need to be redesigned to avoid using insecure time.

Based on this motivation, our publication aims at raising the attention to
possible issues of \acsp{TEE} that may be overlooked, highlighting that not
every \acs{TEE} is the same and each architecture might provide different
capabilities that can influence the threat and attacker model. Thus, when moving
an application into a \acs{TEE}, it is important to carry out a careful security
analysis, making design and implementation changes when necessary.